\subsection{可数序列的使用}

命题 6 是可度量化空间理论中点的可数序列所起作用的根源；对许多问题而言，可以用它们来有效地代替滤子。这是因为可度量化空间中点的邻域滤子（从而还有收敛滤子）由该空间中点的收敛序列确定：因为由于一点的邻域滤子具有可数基，它是比它细的诸基本滤子之交（第一章，§ 6，第 8 号，命题 11），即与收敛到所考察点的诸序列相伴随的基本滤子之交。

另一方面，收敛序列的概念不适用于研究这样的拓扑空间：其中有些点的邻域滤子没有可数基。特别地，可以构造非离散的豪斯多夫拓扑空间，使得在每个点 x 处，x 的邻域的任意可数族之交仍是 \(x (*)\) 的邻域；在这样的空间中，仅有的收敛序列是那些从某个指标起各项都相等的序列。

作为使用可数序列的例子，我们给出下列命题：

命题 8. 在可度量化空间 X 中，点 x 属于 X 的非空子集 A 的闭包，其必要充分条件是存在 A 的一个收敛到 x 的点列。

由第一章，§ 7，第 3 号我们已知该条件是充分的。为看出它是必要的，考察一个可数基本邻域系 \((\mathbf{V}_n)\) （ \(x\) 的），使得 \(\mathbf{V}_{n+1} \subset \mathbf{V}_n\) 对每个 \(n\) 成立。若 \(x\) 属于 A 的闭包，则每个 \(\mathbf{V}_n\) 与 A 相交，且若 \(x_n\) 属于 \(\mathbf{V}_n \cap \mathbf{A}\) ，则序列 \((x_n)\) 收敛到 \(x\) (**)。

由命题 8 我们推出：

命题 9. 度量空间 X 完备的必要充分条件是 X 中每个柯西序列都收敛。

设 \(\hat{\mathbf{X}}\) 是 \(\mathbf{X}\) 的完备化。若存在一点 \(x \in \hat{\mathbf{X}}\) 不属于 \(\mathbf{X}\) ，则存在一个序列 \((x_n)\) （\(\mathbf{X}\) 中的点的序列），它收敛到 \(x\) ，而这是 \(\mathbf{X}\) 中一个不收敛的柯西序列。

命题 10. 设 X 是可度量化空间，f 是 X 到拓扑空间 \(X'\) 中的一个映射。则 f 在点 \(x \in X\) 处连续的必要充分条件是：只要 \((x_n)\) 是 X 中收敛到 x 的点列，序列 \((f(x_n))\) 就收敛到 \(f(x)\) （在 \(X'\) 中）。

该条件是必要的，这由第一章，§ 7，第 4 号，命题 9 的推论 1 得到。为证明它是充分的，考察邻域滤子 \(\mathfrak{B}'\) —— \(f(a)\) 在 \(\mathbf{X}'\) 中的邻域滤子；假设蕴含 \(\overline{f}(\mathfrak{B}')\) 粗于与每个收敛到 \(a\) 的序列相伴随的基本滤子，也就是说粗于每个收敛到 \(a\) 的基本滤子；而这些滤子之交是 \(a\) 的邻域滤子（第一章，§ 6，第 8 号，命题 11）。由此即得结论。

注意，命题 8 与命题 10 在每个点都具有可数基本邻域系的任何空间 X 中都成立。

